\chapter{Martingales and Fair Games}

\begin{center}\itshape
A fair fortune, revealed step by step, is a martingale.\\
Stopping at a bounded time preserves the expectation. Stopping at the first win of one unit does not.
\end{center}

\section{Finite partitions}

Let $\mathcal P=\{A_1,\ldots,A_m\}$ be a partition into events of positive probability, and let $X$ have finite expectation.

\begin{definition}[conditional expectation]\label{def:condexp}
$\mathbb E[X\mid\mathcal P]$ equals $\mathbb E[X\mathbf 1_{A_j}]/P(A_j)$ on the atom $A_j$.
\end{definition}

It is constant on atoms, and $\mathbb E[\mathbb E[X\mid\mathcal P]]=\mathbb E[X]$. If $\mathcal Q$ refines $\mathcal P$, the tower property $\mathbb E[\mathbb E[X\mid\mathcal Q]\mid\mathcal P]=\mathbb E[X\mid\mathcal P]$ is a finite sum: each coarse atom is a union of fine atoms. For square-integrable $X$, it is the orthogonal projection onto the finite-dimensional subspace of variables constant on the atoms, using the inner product $\mathbb E[UV]$ from Volume~1, Chapter~7: the residual has zero mean on each atom, hence is orthogonal to every such step variable.

\section{Martingales and stopping}

\begin{definition}[martingale]\label{def:mart}
A sequence $X_n$, adapted to an increasing sequence of finite partitions, is a martingale when $\mathbb E[X_{n+1}\mid\mathcal P_n]=X_n$.
\end{definition}

A fair random walk is a martingale. Bounded predictable multiples of its steps remain martingale differences.

\begin{theorem}[optional stopping]\label{thm:os}
If $X_n$ is a martingale and $\tau$ is a stopping time with $\tau\le N$ for a fixed integer $N$, then $\mathbb E[X_\tau]=\mathbb E[X_0]$.
\end{theorem}
\begin{proof}
On $\{\tau\ge k\}$, which is determined by $\mathcal P_{k-1}$, the martingale property gives $\mathbb E[X_k\mathbf 1_{\tau\ge k}]=\mathbb E[X_{k-1}\mathbf 1_{\tau\ge k}]$. Write $X_\tau=\sum_{k=0}^N X_k\mathbf 1_{\tau=k}$ and telescope.
\end{proof}

\begin{examplex}[winning one unit]
Let $S_n$ be a fair walk and let $\tau$ be the first time $S_n=1$. Then $S_\tau=1$ and $\mathbb E[S_\tau]=1\ne0$. The time $\tau$ is almost surely finite by Theorem~\ref{thm:polya}, and it is not bounded by any deterministic $N$. Theorem~\ref{thm:os} does not apply.
\end{examplex}

\begin{theorem}[Doob]\label{thm:doob}
If $X_n$ is a martingale with $\mathbb E[X_n^2]<\infty$ and $a>0$, then
\[
 P\Bigl(\max_{k\le n}|X_k|\ge a\Bigr)\le\frac{\mathbb E[X_n^2]}{a^2}.
\]
\end{theorem}
\begin{proof}
Conditional Jensen, or $\mathbb E[Z^2]\ge(\mathbb E Z)^2$ on each atom, shows that $X_k^2$ is a submartingale. Let $\tau$ be the first time among $0,\ldots,n$ when $|X_k|\ge a$, or $n$ if none occurs. Then $a^2\mathbf 1_{\{\max_{k\le n}|X_k|\ge a\}}\le X_\tau^2$. To justify the needed stopped-submartingale inequality without invoking an unstated theorem, write
\[
X_n^2-X_\tau^2=\sum_{k=1}^n\mathbf1_{\{\tau<k\}}(X_k^2-X_{k-1}^2).
\]
Each event $\{\tau<k\}$ belongs to the partition at time $k-1$, and the conditional mean of the corresponding increment is nonnegative. Taking expectations term by term gives $\mathbb E[X_n^2]\ge\mathbb E[X_\tau^2]$. Combined with the displayed indicator bound, this proves the result.
\end{proof}

\begin{theorem}[Azuma]\label{thm:azuma}
If $X_k$ is a martingale adapted to finite partitions and $|X_k-X_{k-1}|\le c_k$ with fixed $c_k\ge0$, then for $t>0$ and $\sum_{k=1}^n c_k^2>0$,
\[
P(X_n-X_0\ge t)\le\exp\bigl(-t^2/(2\sum_{k=1}^n c_k^2)\bigr).
\]
If all $c_k=0$, then $X_n=X_0$ and the probability is $0$ for $t>0$.
\end{theorem}
\begin{proof}
On each atom the increment $D_k=X_k-X_{k-1}$ has conditional mean $0$ and lies in $[-c_k,c_k]$, an interval of length at most $2c_k$. Lemma~\ref{lem:hoeffding}, applied on each atom, gives $\mathbb E[e^{\lambda D_k}\mid\mathcal P_{k-1}]\le\exp(\lambda^2c_k^2/2)$. Multiply by the preceding $\mathcal P_{k-1}$-measurable factor, take expectations, and iterate to obtain
\[
\mathbb E[e^{\lambda(X_n-X_0)}]\le \exp\bigl(\tfrac12\lambda^2\sum_{k=1}^n c_k^2\bigr).
\]
Markov's inequality followed by $\lambda=t/\sum c_k^2$ proves the stated bound when the denominator is positive. If every $c_k=0$, the martingale is constant and the upper-tail probability is $0$ for $t>0$.
\end{proof}

\section{ABRACADABRA}

\begin{theorem}[waiting for a word]\label{thm:monkey}
Independent uniform letters from an alphabet of $26$ symbols produce the word ABRACADABRA at a mean time $26^{11}+26^4+26$.
\end{theorem}
\begin{proof}
A casino accepts one new gambler, with stake $1$, at each integer time. The gambler starting at time $k$ predicts the successive letters of the target word. After $r$ correct guesses their fortune is $26^r$; one wrong guess makes it $0$. After a full successful match the gambler stops. At each step the conditional expected fortune is unchanged, because the next correct letter has probability $1/26$. Let $W_t$ be the sum of all these fortunes after $t$ letters (including any prize completed on that step), and let $M_t=W_t-t$. A new $1$ stake enters at each time, so $M_t$ is a martingale.

The first occurrence time $\tau$ is almost surely finite: in disjoint blocks of eleven letters, the word appears with probability $26^{-11}$ independently, so $P(\tau>11m)\le(1-26^{-11})^m\to0$. For $N\ge1$, put $T_N=\min(\tau,N)$. This is a bounded stopping time, so Theorem~\ref{thm:os} applied to $M$ gives
\[
 \mathbb E[W_{T_N}]=\mathbb E[T_N].
\]
We must justify passing to the unbounded time rather than simply invoking optional stopping there. At any time before the first occurrence, a surviving gambler must have started among the last eleven letters: an older gambler either lost or completed the target, which would already have stopped the clock. At $\tau$ the same bound holds, including the completed gambler. Each fortune is at most $26^{11}$, so $0\le W_{T_N}\le 11\cdot26^{11}$ for every $N$. As $N\to\infty$, $T_N\uparrow\tau$ and $W_{T_N}\to W_\tau$ almost surely; bounded convergence on the wealth and monotone convergence on the time give $\mathbb E[\tau]=\mathbb E[W_\tau]<\infty$.

At the first completed ABRACADABRA the surviving gamblers are exactly those matching suffixes that are also prefixes. The full eleven-letter word pays $26^{11}$, the suffix ABRA of length four pays $26^4$, and the suffix A of length one pays $26$. There are no other proper prefix--suffix matches, so $W_\tau=26^{11}+26^4+26$ deterministically. Therefore
\[
 \mathbb E[\tau]=26^{11}+26^4+26.
\]
The stopping step is valid because we proved the bound and passed through $T_N$, not because unbounded optional stopping holds in general.
\end{proof}

A pattern with no prefix equal to a proper suffix has mean waiting time equal to the single leading power. Chapter~12 constructs a continuous martingale with Gaussian increments.

\section{Exercises}

\begin{exercise}[$\star$]\label{ex:tower}
Prove the tower property by averaging fine atoms inside a coarse atom.
\end{exercise}

\begin{exercise}[$\star\star$]\label{ex:stopmean}
Let $\tau$ be the first time a fair walk hits $2$, or time $10$ if it has not. Show $\mathbb E[S_\tau]=0$, and that $S_\tau$ is not almost surely $2$.
\end{exercise}

\begin{exercise}[$\star\star$]\label{ex:azumacoin}
Apply Theorem~\ref{thm:azuma} to the fair walk and compare the bound with Theorem~\ref{thm:hoeffding}.
\end{exercise}

\begin{exercise}[$\star\star$]\label{ex:pattern}
For a pattern of length $\ell$ on an alphabet of size $q$, with no proper prefix equal to a suffix, show that the mean waiting time is $q^\ell$.
\end{exercise}

\section{Hints}

\begin{hint}{ex:tower}
A coarse atom is a union of fine atoms.
\end{hint}
\begin{hint}{ex:stopmean}
The time is at most $10$.
\end{hint}
\begin{hint}{ex:azumacoin}
Steps change the walk by $1$.
\end{hint}
\begin{hint}{ex:pattern}
Only one gambler is still winning at the end.
\end{hint}

\section{Solutions}

\begin{solution}{ex:tower}
\writeup On a coarse atom the fine conditional expectation is constant on each fine piece. Averaging those constants with their probabilities reproduces the coarse average.
\end{solution}

\begin{solution}{ex:stopmean}
\writeup Theorem~\ref{thm:os} applies because $\tau\le 10$, so the expectation is $0$. The walk can stay inside $\{-1,0,1\}$ for ten steps, an event of positive probability, and then $S_\tau=S_{10}\ne2$.
\end{solution}

\begin{solution}{ex:azumacoin}
\writeup The differences are bounded by $1$, so Theorem~\ref{thm:azuma} gives $\exp(-t^2/(2n))$. Theorem~\ref{thm:hoeffding} for summands in an interval of length $2$ gives the same exponent, because $\sum(b_i-a_i)^2=4n$ and $-2t^2/(4n)=-t^2/(2n)$.
\end{solution}

\begin{solution}{ex:pattern}
\writeup Restart the casino construction of Theorem~\ref{thm:monkey} with an alphabet of size $q$ and a target of length $\ell$. Disjoint length-$\ell$ blocks have a fixed positive chance $q^{-\ell}$ to match, so the waiting time is finite almost surely. Stop the total-wealth-minus-number-of-stakes martingale at $\tau\wedge N$; before the first match at most $\ell$ gamblers can remain active, each worth at most $q^\ell$. Their total wealth is uniformly bounded, so bounded convergence on wealth and monotone convergence on time justify passing $N\to\infty$. At the first completion only the gambler who started $\ell$ letters earlier still holds wealth: a different winning gambler would require a proper suffix equal to a prefix, contrary to the hypothesis. Its fortune is $q^\ell$, hence the mean waiting time is $q^\ell$.
\end{solution}

\section*{Sources and further reading}

\begin{enumerate}
\item J.~L.~Doob, \emph{Stochastic Processes}, Wiley, 1953. Theorem~\ref{thm:doob} in the discrete form used here.
\item D.~Williams, \emph{Probability with Martingales}, Cambridge University Press, 1991. The casino proof of Theorem~\ref{thm:monkey}.
\end{enumerate}
